\documentclass[../../main.tex]{subfiles}
\begin{document}
% SOURCE p=235; printed=220; object=ch07-regularity-exercises-title; region=top
\Needspace{12\baselineskip}
\subsection*{习题}
% SOURCE p=235; printed=220; object=ch07-ex07; number=7; region=top
\begin{exercise}\label{exercise:ch07-restricted-radon}
若 $\mu$ 为 $X$ 上的 $\sigma$ 有限 Radon 测度，$A\in\mathscr B_X$，则由 $\mu_A(E)=\mu(E\cap A)$ 定义的 Borel 测度 $\mu_A$ 是 Radon 测度。（另见\cref{exercise:ch07-zero-weight-failure}。）
\end{exercise}
% SUPPLEMENT: id=v14-ch07-ex07-solution; parent=ch07-ex07; kind=solution+motivation+hypotheses; origin=independent-derivation
\begin{solution}
\textbf{截取可测集上的质量（重构新增）。}交 $A$ 保持不交可数并，故 $\mu_A$ 可数可加；且 $\mu_A(K)\leqslant\mu(K)<\infty$ 对所有紧集成立。对任意 Borel 集 $E$，$E\cap A$ 为 $\sigma$ 有限，由\cref{proposition:ch07-sigma-finite-inner}，
\begin{align*}
\mu_A(E)=\mu(E\cap A)=\sup_{K\subseteq E\cap A,\ K\text{ 紧}}\mu(K)
\leqslant\sup_{K\subseteq E,\ K\text{ 紧}}\mu_A(K)\leqslant\mu_A(E).
\end{align*}
因此 $\mu_A$ 在所有 Borel 集上内正则，特别在开集上如此。若 $\mu_A(E)=\infty$，外正则自动成立；否则，由\subcref{proposition:ch07-open-closed-approximation}{item:ch07-open-closed-epsilon}选开集 $U\supseteq E$ 使 $\mu(U\setminus E)<\varepsilon$，于是
\begin{align*}
\mu_A(U)=\mu_A(E)+\mu_A(U\setminus E)
\leqslant\mu_A(E)+\mu(U\setminus E)<\mu_A(E)+\varepsilon.
\end{align*}
故 $\mu_A$ 外正则，三项 Radon 条件均已证明。
\end{solution}
\begin{remark}
\textbf{动机与逐条件说明（重构新增）。}$A$ 只需 Borel，不必开或闭；其作用是把 $A^c$ 上的质量去掉。$\sigma$ 有限性使 $E\cap A$ 可内逼近，也保证无论原 $\mu(E)$ 是否有限都能控制开外罩的\emph{差集}质量；只用 $\mu_A(E)<\infty$ 不能推出 $\mu(E)<\infty$。后面的反例表明缺少全局 $\sigma$ 有限性时，截取甚至一个开集也可能破坏外正则性。
\end{remark}
% END SUPPLEMENT: id=v14-ch07-ex07-solution

% SOURCE p=235; printed=220; object=ch07-ex08; number=8; region=upper
\begin{exercise}\label{exercise:ch07-integrable-weight-radon}
设 $\mu$ 为 $X$ 上的 Radon 测度。若 $\phi\in L^1(\mu)$ 且 $\phi\geqslant0$，则 $\nu(E)=\int_E\phi\,\mathrm d\mu$ 为 Radon 测度。（使用\cref{corollary:ch03-integral-absolute-continuity}。）
\end{exercise}
% SUPPLEMENT: id=v14-ch07-ex08-solution; parent=ch07-ex08; kind=solution+motivation+hypotheses; origin=independent-derivation
\begin{solution}
\textbf{有限加权质量的紧逼近（重构新增）。}非负积分的单调收敛保证 $\nu$ 是 Borel 测度，且 $\nu(X)=\|\phi\|_1<\infty$。固定 Borel 集 $E$ 及 $\varepsilon>0$，令 $E_n=E\cap\{\phi\geqslant1/n\}$。有 $\mu(E_n)\leqslant n\nu(X)<\infty$，且 $E_n\uparrow E\cap\{\phi>0\}$；由 $\nu(E\cap\{\phi=0\})=0$ 及 $\nu(E)<\infty$，可取 $n$ 使 $\nu(E\setminus E_n)<\varepsilon/2$。

积分的绝对连续性给出 $\delta>0$，使 $\mu(H)<\delta$ 时 $\int_H\phi\,\mathrm d\mu<\varepsilon/2$。由\cref{proposition:ch07-sigma-finite-inner}，取紧集 $K\subseteq E_n$ 使 $\mu(E_n\setminus K)<\delta$。于是 $\nu(E\setminus K)<\varepsilon$，证明在所有 Borel 集上内正则。再对 $E^c$ 应用此结论，取紧集 $L\subseteq E^c$ 使 $\nu(E^c\setminus L)<\varepsilon$。$U=X\setminus L$ 开、包含 $E$，且 $\nu(U\setminus E)<\varepsilon$，所以外正则。紧集有限由总质量有限直接得到，故 $\nu$ 为 Radon 测度。
\end{solution}
\begin{remark}
\textbf{动机与逐条件说明（重构新增）。}非负密度保证所得对象是非负测度；$L^1$ 条件既给出有限总质量，又给出小 $\mu$ 测度集上的小积分，二者均在证明中使用。$\mu$ 本身可以不 $\sigma$ 有限；只需先在 $\phi$ 的正水平集上局部化。$L^1$ 元若按等价类理解，选非负 Borel 代表即可，改变零集上的值不改变 $\nu$。由紧内逼近取补得到外逼近时，有限的\emph{新}总质量使尾误差可控。
\end{remark}
% END SUPPLEMENT: id=v14-ch07-ex08-solution

% SOURCE p=235; printed=220; object=ch07-ex09; number=9; region=upper-middle
\begin{exercise}\label{exercise:ch07-positive-continuous-weight}
设 $\mu$ 为 $X$ 上的 Radon 测度，$\phi\in C(X,(0,\infty))$。令 $\nu(E)=\int_E\phi\,\mathrm d\mu$，并令 $\nu'$ 为 $C_c(X)$ 上泛函 $f\mapsto\int f\phi\,\mathrm d\mu$ 对应的 Radon 测度。
\begin{enumerate}[label=\alph*.,ref=(\alph*)]
\item\label{item:ch07-positive-weight-open} 若 $U$ 开，则 $\nu(U)=\nu'(U)$。（对 $\phi\chi_U$ 应用\cref{corollary:ch07-lsc-integral-minorants}。）
\item\label{item:ch07-positive-weight-outer} $\nu$ 在所有 Borel 集上外正则。（提示：开集 $V_k=\{x:2^k<\phi(x)<2^{k+2}\}$，$k\in\mathbb Z$，覆盖 $X$。）
\item\label{item:ch07-positive-weight-radon} $\nu=\nu'$，因而 $\nu$ 为 Radon 测度。（另见\cref{exercise:ch07-zero-weight-failure}。）
\end{enumerate}
\end{exercise}
% SUPPLEMENT: id=v14-ch07-ex09-solution; parent=ch07-ex09; kind=solution+motivation+hypotheses; origin=independent-derivation
\begin{solution}
\textbf{正连续密度与层带外逼近（重构新增）。}$\phi$ 连续且在每个紧集上有界，所以 $f\phi\in C_c(X)$ 对每个 $f\in C_c(X)$ 成立；积分定义正线性泛函，$\nu'$ 因而存在。$\nu$ 可数可加且在紧集上有限。

(a) $\phi\chi_U$ 非负下半连续：对 $a<0$ 严格上水平集为 $X$，对 $a\geqslant0$ 为 $U\cap\{\phi>a\}$，均开。由所提示推论，
\begin{align*}
\nu(U)=\sup\left\{\int g\,\mathrm d\mu:g\in C_c(X),\ 0\leqslant g\leqslant\phi\chi_U\right\}.
\end{align*}
对每个候选 $g$，$h=g/\phi$ 连续且紧支，并且 $0\leqslant h\leqslant\chi_U$。于是 $\int g\,\mathrm d\mu=\int h\phi\,\mathrm d\mu=\int h\,\mathrm d\nu'\leqslant\nu'(U)$。取上确界得 $\nu(U)\leqslant\nu'(U)$。反过来，由 Riesz 的开集公式，
\begin{align*}
\nu'(U)=\sup_{f\prec U}\int f\phi\,\mathrm d\mu\leqslant\int_U\phi\,\mathrm d\mu=\nu(U).
\end{align*}

(b) 若 $\nu(E)=\infty$，外正则自动成立。若有限，给定 $\varepsilon>0$，选正数 $\delta_k$（$k\in\mathbb Z$）使 $\sum_k\delta_k<\varepsilon$。令 $E_k=E\cap V_k$，则 $\mu(E_k)\leqslant2^{-k}\nu(E)<\infty$。由 $\mu$ 的外正则性，取开集 $O_k\supseteq E_k$ 使 $\mu(O_k\setminus E_k)<\delta_k2^{-k-2}$，令 $U_k=O_k\cap V_k$。在 $V_k$ 上 $\phi<2^{k+2}$，且 $U_k\setminus E\subseteq O_k\setminus E_k$，所以 $\nu(U_k\setminus E)<\delta_k$。$U=\bigcup_{k\in\mathbb Z}U_k$ 开且包含 $E$，于是
\begin{align*}
\nu(U\setminus E)\leqslant\sum_k\nu(U_k\setminus E)<\varepsilon,
\end{align*}
证明外正则性。

(c) 两测度均在 Borel 集上外正则，并在开集上相等，因此对每个 Borel 集 $E$，
\begin{align*}
\nu(E)=\inf_{U\supseteq E,\ U\text{ 开}}\nu(U)
=\inf_{U\supseteq E,\ U\text{ 开}}\nu'(U)=\nu'(E).
\end{align*}
故 $\nu$ 为 Radon 测度。
\end{solution}
\begin{remark}
\textbf{动机与逐条件说明（重构新增）。}严格正性使除以 $\phi$ 合法，并使可数个双边有界层带 $V_k$ 覆盖整个空间；不要求全局正下界或全局上界。连续性保证这些层带开、商函数连续、紧集上的密度有界。$h\leqslant\chi_U$ 不必意味着 $h\prec U$，因为支集可能触到边界；(a) 只使用积分比较，从未误用这个更强条件。每层的正下界把有限 $\nu$ 质量变为有限 $\mu$ 质量，上界把 $\mu$ 误差变为 $\nu$ 误差，可求和预算再拼回全空间。这替代了上一题的全局可积性条件。
\end{remark}
% END SUPPLEMENT: id=v14-ch07-ex09-solution

% SOURCE p=235; printed=220; object=ch07-ex10; number=10; region=middle
\begin{exercise}\label{exercise:ch07-real-lone-semicontinuous}
若 $\mu$ 为 Radon 测度，实值 $f\in L^1(\mu)$，则对每个 $\varepsilon>0$，存在下半连续函数 $g$ 和上半连续函数 $h$，使 $h\leqslant f\leqslant g$ 且 $\int(g-h)\,\mathrm d\mu<\varepsilon$。
\end{exercise}
% SUPPLEMENT: id=v14-ch07-ex10-solution; parent=ch07-ex10; kind=solution+motivation+hypotheses; origin=independent-derivation
\begin{solution}
\textbf{正负部分的交叉上下包（重构新增）。}$f^+,f^-$ 非负可积，它们的正支撑均 $\sigma$ 有限。由\cref{proposition:ch07-semicontinuous-integral-approximation}及其证明，可选非负下半连续函数 $u_\pm\geqslant f^\pm$，以及\emph{有限值}非负上半连续函数 $0\leqslant v_\pm\leqslant f^\pm$，使
\begin{align*}
\int(u_\pm-f^\pm)\,\mathrm d\mu<\varepsilon/4,\qquad
\int(f^\pm-v_\pm)\,\mathrm d\mu<\varepsilon/4.
\end{align*}
有限值下包可直接采用该命题证明中的有限个紧集示性函数之和；若相应积分为零，取零函数。令
\begin{align*}
g=u_+-v_-,\qquad h=v_+-u_-.
\end{align*}
$-v_-$ 为有限值下半连续函数，所以 $g$ 下半连续，允许 $+\infty$；$-u_-$ 上半连续，$v_+$ 有限，故 $h$ 上半连续，允许 $-\infty$。两式都没有 $\infty-\infty$。由各点的比较，$h\leqslant f^+-f^-=f\leqslant g$。又有非负函数恒等式
\begin{align*}
g-h=(u_+-f^+)+(f^+-v_+)+(u_--f^-)+(f^--v_-).
\end{align*}
积分后四个误差相加严格小于 $\varepsilon$，得到结论。
\end{solution}
\begin{remark}
\textbf{动机与逐条件说明（重构新增）。}上下包必须交叉组合：上界用正部分的上包减去负部分的下包，下界反之。实值条件使逐点顺序和正负分解有意义；可积性确保两部分的积分有限及其正支撑 $\sigma$ 有限。有限值下包的选择避免在有无穷值的点出现未定义运算。结论按 Borel 实值代表逐点陈述；对完备化等价类若只给出类，应先选 Borel 代表，此时类之间的比较自然按几乎处处解释。
\end{remark}
% END SUPPLEMENT: id=v14-ch07-ex10-solution

% SOURCE p=235; printed=220; object=ch07-ex11; number=11; region=middle
\begin{exercise}\label{exercise:ch07-atomless-prescribed-mass}
设 $\mu$ 为 $X$ 上的 Radon 测度，且对所有 $x\in X$ 有 $\mu(\{x\})=0$。若 $A\in\mathscr B_X$ 满足 $0<\mu(A)<\infty$，则任给 $0<\alpha<\mu(A)$，存在 Borel 集 $B\subseteq A$ 使 $\mu(B)=\alpha$。
\end{exercise}
% SUPPLEMENT: id=v14-ch07-ex11-solution; parent=ch07-ex11; kind=solution+motivation+hypotheses; origin=independent-derivation
\begin{solution}
\textbf{有限小块分割与递增填充（重构新增）。}先证明：对任意有限测度 Borel 集 $E$ 及 $\delta>0$，存在 $E$ 的一个有限 Borel 分割，使每片测度小于 $\delta$。由内正则取紧集 $K\subseteq E$，使 $\mu(E\setminus K)<\delta$。对每个 $x\in K$，由单点零质量和外正则性，取开邻域 $U_x$ 使 $\mu(U_x)<\delta$。有限多个 $U_{x_i}$ 覆盖 $K$；把 $K\cap U_{x_i}$ 依次去掉前面的部分得到不交的 Borel 片，每片包含于相应 $U_{x_i}$，测度小于 $\delta$。再加入 $E\setminus K$，便得所需分割（空片可删）。

现在递归构造 $C_n\subseteq A$，使它们递增、$\mu(C_n)\leqslant\alpha$，并且 $\alpha-\mu(C_n)<2^{-n}$。取 $C_0=\varnothing$。若某一步 $\mu(C_{n-1})=\alpha$，此后保持该集即可。否则令 $d=\alpha-\mu(C_{n-1})>0$、$R=A\setminus C_{n-1}$。因为 $\alpha<\mu(A)$，有 $\mu(R)>d$。把 $R$ 按上段分成有限个测度小于 $2^{-n}$ 的 Borel 片 $R_1,\ldots,R_m$。其部分和最终大于 $d$，取第一个使 $\sum_{j=1}^k\mu(R_j)>d$ 的指标 $k$；令 $D=\bigcup_{j<k}R_j$。于是
\begin{align*}
0\leqslant d-\mu(D)<\mu(R_k)<2^{-n}.
\end{align*}
取 $C_n=C_{n-1}\cup D$，满足所需条件。最后令 $B=\bigcup_{n=1}^\infty C_n$。它为 $A$ 内的 Borel 集，由从下连续性，$\mu(B)=\lim_n\mu(C_n)=\alpha$。
\end{solution}
\begin{remark}
\textbf{动机与逐条件说明（重构新增）。}单点无质量与外正则性先产生任意小质量的邻域，紧内逼近把这些邻域化为有限分割；然后用逐次填充精确取得目标质量。不能只声称“存在任意小的正质量集”而省略填充步骤。有限的 $\mu(A)$ 使内逼近留下小尾部，且所有差集的测度相减合法；$0<\alpha<\mu(A)$ 保证每轮残余质量严格超过待填差额。端点 $\alpha=0,\mu(A)$ 分别可取空集和 $A$，但原题关注非平凡情形。若存在正质量原子，则一般不能把它任意细分，例如单位点质量不能取得 $1/2$。
\end{remark}
% END SUPPLEMENT: id=v14-ch07-ex11-solution
\subfile{section-02-counterexamples.tex}
\end{document}
